% =====================================================================
%  MOSAIC — 사용 수식 정리 (이해용 레퍼런스)
%  컴파일: xelatex formulas.tex   (권장; 한글)
%          또는 pdflatex + kotex
%  내용은 pfire/ 소스의 실제 구현과 1:1 대응한다(모듈명을 각주로 표기).
% =====================================================================
\documentclass[11pt,a4paper]{article}

\usepackage{kotex}
\usepackage[margin=2.2cm]{geometry}
\usepackage{amsmath,amssymb}
\usepackage{booktabs}
\usepackage{enumitem}
\usepackage[hidelinks]{hyperref}
\usepackage{xcolor}

\newcommand{\src}[1]{{\footnotesize\color{gray}\texttt{#1}}}
\newcommand{\R}{\mathbb{R}}
\setlength{\parindent}{0pt}
\setlength{\parskip}{4pt}

\title{\textbf{MOSAIC 산불위험 조기경보 — 사용 수식 정리}\\[4pt]
\large 전력설비(전주) 산불위험 모델에 쓰인 식과 그 의미}
\author{이해용 레퍼런스 (코드 \texttt{pfire/} 구현과 1:1 대응)}
\date{2026-06-20}

\begin{document}
\maketitle

\begin{abstract}
\noindent
강원 전주 \textbf{138만 개}에 산불 위험/안전$\{0,1\}$을 매기는 문제다. 정답 라벨이
없으므로(presence-only) 과거 발화점은 \emph{학습 타깃이 아니라 검증·임계값 앵커}로만
쓴다. 핵심 아이디어는 위험을 \emph{발화}$\times$\emph{확산취약}$\times$\emph{기상}의
곱으로 분해하고, 각 항을 물리식/지역 혼합전문가로 추정한 뒤, 불확실성(베이지안 사후)과
커버리지(conformal)를 붙이는 것이다. 이 문서는 그 모든 식을 모듈 순서대로 모아
\emph{무엇을 왜 그렇게 계산했는지} 한 곳에서 본다.
\end{abstract}

\tableofcontents
\vspace{6pt}\hrule

% =====================================================================
\section{기호 약속}
% =====================================================================
\begin{tabular}{ll}
\toprule
기호 & 의미 \\
\midrule
$p$ & 전주 인덱스 ($N=1{,}387{,}831$) \\
$t$ & 날짜(산불조심기간 2$\sim$5월) \\
$r$ & 체제(regime): 영동/영서/산간 \\
$g(p)$ & 전주 $p$ 가 속한 기상 격자($\sim$13\,km) \\
$I,S,W\in[0,1]$ & 발화 성향 / 확산·노출 취약 / 기상 위험 \\
$h,R\in[0,1]$ & 일 위험 / 시즌 위험 \\
$y_g,\,n_g$ & 지역 $g$ 의 발화수(앵커) / 전주수(노출) \\
\bottomrule
\end{tabular}

모든 성분은 $[0,1]$ 로 정규화하고 NaN/범위를 강제 검증한다(\src{hazard.\_validate\_unit}).

% =====================================================================
\section{위험의 곱셈 분해}
% =====================================================================
\src{pfire/hazard.py}

전주 하루치 위험을 세 항의 곱으로 본다:
\begin{equation}
h(p,t)\;=\;\underbrace{I(p)}_{\text{발화 성향}}\;\times\;
\underbrace{S(p)}_{\text{확산·노출 취약}}\;\times\;
\underbrace{W(g(p),t)}_{\text{그날 기상}}.
\end{equation}

\textbf{왜 곱인가:} 셋 중 하나라도 0이면 위험 0이라는 \emph{AND} 성격이 물리적으로
타당하고(불씨가 없거나/탈 게 없거나/날이 궂지 않으면 안 난다), 각 항을 따로 해석해
``왜 위험한가''를 항목별로 설명할 수 있다.

\textbf{시즌 집계 $R(p)$.} 산불조심기간을 모아 하나의 시즌 위험으로 집계한다.
\begin{align}
\text{(MVP, 단일 } W)\quad & R(p)=I(p)\,S(p)\,W(p),\\
\text{(일별 } T\text{일)}\quad & R(p)=\exp\!\Big(\tfrac{1}{T}\textstyle\sum_{t=1}^{T}
\log\big[\mathrm{clip}(I\,S\,W_t,\,\varepsilon,1)\big]\Big).
\end{align}
일별 경로는 \emph{기하평균}(로그합 평균)이라 단발 고위험일 하나가 시즌 위험을 폭주시키는
것을 완화한다.

% =====================================================================
\section{발화 성향 $I(p)$ — 지역 혼합전문가(MoE)}
% =====================================================================
\src{pfire/experts.py, pfire/regimes.py}

영동/영서/산간은 발화 메커니즘이 다르다(영동=송전선·양간풍, 영서=도로·입산자,
산간=연료·산림). 전 지역 단일 모델은 평균에 묻히고, 지역별 개별 모델은 데이터가 쪼개져
과적합한다. 그 중간(부분 풀링)을 \emph{soft 게이트 혼합전문가}로 구현한다.

\subsection{피처 정규화 (모두 $[0,1]$, 클수록 위험)}
거리 피처는 \emph{가까울수록 위험}이므로 지수 감쇠로 정규화한다:
\begin{equation}
\mathrm{feat}_{\text{dist}}(p)=\exp\!\big(-\,d(p)/\text{scale}\big),\qquad
\text{scale}=\begin{cases}
100\,\text{m} & \text{산림}\\
300\,\text{m} & \text{도로}\\
1500\,\text{m} & \text{송전선}\\
3000\,\text{m} & \text{변전소(선택)}
\end{cases}
\end{equation}
나머지: FWI·양간일은 min--max 정규화, 연료는 $\mathrm{clip}(\mu_{\text{flam}},0,1)$,
토지피복은 도메인 사전치 $\mathrm{lc\_ignition}\in[0,1]$ 조회(산림 1.00, 초지 0.25,
$\dots$, 수역 0.00).

\subsection{체제 전문가 (같은 피처, 다른 가중)}
체제 $r$ 전문가는 피처들의 \emph{가중평균}이다(가중합을 가중치합으로 나눠 $[0,1]$ 보장):
\begin{equation}
E_r(p)\;=\;\frac{\sum_{k} w_{r,k}\,\mathrm{feat}_k(p)}{\sum_{k} w_{r,k}},\qquad
k\in\{\text{forest, road, powerline, fwi, yanggan, fuel, landcover}\}.
\end{equation}
가중 $w_{r,k}$ 는 공간블록 CV 발화점 recall 을 목적으로 LHS+좌표상승 튜닝했다
(\src{config.EXPERT\_WEIGHTS}; 예: 영동은 forest 0.33·fuel 0.33, 영서는
powerline 0.35·forest 0.35).

\subsection{soft 게이트 (체제 가중)}
게이트 피처 $z_p=(\text{경도, 고도, 양간일})$ 를 z-표준화하고, 체제 앵커 $\mu_r$
(시군 매핑된 전주들의 표준화 피처 중심)과의 거리로 softmax 한다:
\begin{equation}
g_r(p)\;=\;\frac{\exp\!\big(-\lVert z_p-\mu_r\rVert^2/\tau\big)}
{\sum_{r'}\exp\!\big(-\lVert z_p-\mu_{r'}\rVert^2/\tau\big)},
\qquad \textstyle\sum_r g_r(p)=1,\quad \tau=0.6.
\end{equation}
경계 전주는 자연 블렌딩되고, 소표본 시군도 피처로 안정 배치된다(하드 분할이 아님).

\subsection{MoE 결합}
\begin{equation}
\boxed{\,I(p)\;=\;\sum_{r} g_r(p)\,E_r(p)\,}\;\in[0,1].
\end{equation}

% =====================================================================
\section{기상 위험 $W$}
% =====================================================================
\src{pfire/weather.py}

\subsection{시즌 통계 성분}
전주 사전계산 통계를 결합(각 min--max 정규화):
\begin{equation}
W_{\text{season}}(p)=0.5\,\underbrace{\mathrm{fwi\_q90}}_{\text{강도}}
+0.3\,\underbrace{\mathrm{hdd}}_{\text{고위험일 빈도}}
+0.2\,\underbrace{\mathrm{yanggan}}_{\text{강풍}}.
\end{equation}

\subsection{일별 ISI/FFMC 동역학 성분}
관측소별 산불조심기간 일별 분포를 시즌 집계한 뒤 최근접 관측소로 전주에 부여한다.
\begin{align}
\text{ISI 성분} &= 0.5\,(\text{고-ISI일수비율}) + 0.25\,(\overline{\text{ISI}})
+ 0.25\,(\text{ISI}_{q90}),\\
\text{FFMC 성분} &= \text{고-FFMC일수비율},\qquad
\text{양간 성분}=\mathrm{yanggan\_days},\\
W_{\text{daily}}(p) &= 0.05\,(\text{ISI 성분}) + 0.75\,(\text{FFMC 성분})
+ 0.20\,(\text{양간 성분}).
\end{align}
고-ISI 임계 $10.0$, 고-FFMC 임계 $90.0$(관측소 일별 $q90$ 정합). 발화점 앵커가
인간발화 위주라 \emph{발화준비(FFMC)}가 본체, \emph{확산준비(ISI)}는 소수 유지.

\subsection{블렌드(기본 $W$)}
일별만 쓰면 전주를 관측소 시즌집계로 매핑해 \emph{전주고유 극값}(예: 2019 고성 인근)을
희석한다. 둘을 선형 결합해 극값 보존과 recall 이득을 함께 취한다:
\begin{equation}
\boxed{\,W(p)\;=\;w_s\,W_{\text{season}}(p)+(1-w_s)\,W_{\text{daily}}(p)\,},
\qquad w_s=0.82.
\end{equation}

% =====================================================================
\section{확산·노출 취약 $S(p)$}
% =====================================================================
\src{pfire/hazard.py (blend\_exposure*)}

정적 $S_{\text{static}}$(사전계산)을 풍하 노출과 \emph{전주별 조건부}로 블렌딩한다.
전역 균일 희석을 막기 위해 영동·고양간 전주에서만 블렌드 강도를 키운다:
\begin{align}
\alpha_p &= \alpha_{\text{base}}\cdot \underbrace{g_{\text{영동}}(p)}_{\text{영동 게이트}}
\cdot \underbrace{\widehat{\mathrm{yanggan}}(p)}_{\text{양간강도 백분위}}\ \in[0,\alpha_{\text{base}}],\\
S(p) &= (1-\alpha_p)\,S_{\text{static}}(p) + \alpha_p\,\widetilde{P}_{\text{expo}}(p),
\end{align}
여기서 $\widetilde{P}_{\text{expo}}$ 는 노출확률의 순위 백분위(0 다수 분포 강건화).

% =====================================================================
\section{풍하 노출 — 두 가지 커널}
% =====================================================================

\subsection{v1: OR 확산 커널 (몬테카를로)}
\src{pfire/exposure.py, exposure\_engine.py}

발화후보(=$I$ 상위 2\% 전주)에서 바람을 표집해 비등방 확산을 시뮬레이션한다. 시뮬 $s$
마다 풍하방향 $\theta=\text{풍향}+180^\circ$, 정렬도 $\mathrm{align}=\cos(\angle(\text{방위},\theta))$:
\begin{align}
L &= L_0\,\big(1+\alpha\,\max(0,\mathrm{align})\,\tfrac{ws}{5}\big),\quad L_0{=}0.8\text{km},\ \alpha{=}1.5,\\
\Pr(\text{도달})&=\exp(-d/L)\cdot \mathrm{fuel}\cdot(1+\beta\,\mathrm{southness}),\quad \beta{=}0.3,\\
P_{\text{expo}}(p)&=\frac{1}{S}\textstyle\sum_{s}\mathbf{1}[\text{어느 발화원이라도 } p \text{ 도달}].
\end{align}
\textbf{문제:} 발화원이 조밀한 영동에서 OR 확률이 $P\approx1$ 로 \emph{포화}되어 변별력을
잃는다 $\Rightarrow$ 아래 v2.2 로 재설계.

\subsection{v2.2: 변별 가능한 발화가중 기대도즈}
\src{pfire/exposure\_v2.py}

OR(포화) 대신 \emph{누적합}(비포화)으로 바꾸고, Anderson 타원 커널과 격자 국지정규화를
더한다. 풍속 $\to$ 길이/폭비(LWR):
\begin{equation}
U=2.237\,ws\ (\text{mph}),\quad
\mathrm{LWR}(ws)=\min\!\big(0.936\,e^{0.2566U}+0.461\,e^{-0.1548U}-0.397,\ 8\big).
\end{equation}
발화원 $g$ 에 대해 평행/수직거리 $(d_\parallel,d_\perp)$ 로 타원 반경 $r$ 을 만들고,
풍하($d_\parallel\!\ge\!0$)는 길게·풍상은 짧게 둔다:
\begin{align}
L_\parallel &= \begin{cases}\mathrm{LWR}\cdot W_\perp & d_\parallel\ge 0\ (\text{풍하})\\ L_{\text{back}} & d_\parallel<0\ (\text{풍상})\end{cases},
\qquad W_\perp=0.35,\ L_{\text{back}}=0.35,\\
r &= \sqrt{(d_\parallel/L_\parallel)^2+(d_\perp/W_\perp)^2},\\[2pt]
\mathrm{dose}(p) &= \sum_{g} I_g\ \Big\langle e^{-r}\Big\rangle_{\text{바람}}\ \mathrm{fuel}_p\,(1+\beta\,s_p).
\end{align}
체제 결합 $\mathrm{dose}=\sum_r g_r\,\mathrm{dose}_r$ 후 \emph{격자 국지정규화}로
``영동은 다 높음'' 배경을 제거:
\begin{equation}
\widetilde{e}(p)=\frac{\mathrm{dose}(p)}{\overline{\mathrm{dose}}_{\,\mathrm{cell}(p)}+\varepsilon},
\qquad \mathrm{dose01}(p)=\frac{\mathrm{percentile}\big(\widetilde{e}(p)\big)}{100}.
\end{equation}

\subsection{결정에 결합 (확률 OR)}
발화 랭킹을 보존하면서 풍하 전주를 상대적으로 끌어올린다:
\begin{equation}
\boxed{\,R'(p)=1-\big(1-R(p)\big)\big(1-w\,\mathrm{dose01}(p)\big)\,},\qquad w=0.25\ (\text{보수}).
\end{equation}
검증: 영동 변별 IQR $0.166\!\to\!0.290$, 공간CV recall $+0.0071$(5 fold-seed 신호).

% =====================================================================
\section{지역율 계층 보정 — 경험적 베이즈(EB)}
% =====================================================================
\src{pfire/hierarchy.py}

발화점이 희소($\sim$700개)하고 일부 시군은 발화 0$\sim$소수다. 관측율을 그대로 쓰면
극단 추정으로 과적합하므로, Poisson--Gamma EB 축소로 부모 단계 평균으로 끌어당긴다
(전역$\to$체제$\to$시군$\to$격자, top-down).

지역 $g$ 의 관측율 $\hat\lambda_g=y_g/n_g$. 부모율 $m$ 과 형제율 분산 $v$ 로 Gamma 사전을
적률추정($\alpha=m^2/v,\ \beta=m/v$)하면 사후평균은
\begin{equation}
\tilde\lambda_g=\frac{y_g+\alpha}{n_g+\beta}
=w_g\,\hat\lambda_g+(1-w_g)\,m,\qquad w_g=\frac{n_g}{n_g+\beta}.
\end{equation}
표본 큰 지역은 관측율에 가깝고($w_g\to1$), 0발화·소표본은 부모율 $m$ 으로 축소된다
($w_g\to0$). 전주 배율은 격자 EB율을 전역 EB율로 나눈 \emph{상대배율}(전역평균$\approx$1).

% =====================================================================
\section{불확실성 — 베이지안 사후 + MC 전파}
% =====================================================================
\src{pfire/posterior.py}

\subsection{Poisson--Gamma 켤레 사후 (닫힌형)}
지역 발화수에 켤레 사전을 두면 사후가 닫힌형으로 나온다(MCMC 불필요):
\begin{equation}
\lambda_g\,\big|\,y_g\ \sim\ \mathrm{Gamma}\!\big(\alpha_0+y_g,\ \beta_0+n_g\big),
\quad \mathbb{E}=\tfrac{\alpha}{\beta},\ \mathrm{Var}=\tfrac{\alpha}{\beta^2}.
\end{equation}
zero-event 지역은 $\alpha\!\to\!\alpha_0$(부모 수렴)이고 노출 작아 $\beta$ 도 작으므로
\emph{사후 분산이 커진다}(불확실성 자동 표현).

\subsection{위험 사후예측(몬테카를로)}
각 draw $r$ 에서 지역율 사후배율 $m^{(r)}$(전역평균$\approx$1)을 물리 base risk 에 곱한다:
\begin{equation}
R^{(r)}(p)=\mathrm{clip}\big(R_{\text{base}}(p)\cdot m^{(r)}(p),\,0,1\big),
\end{equation}
전주별 평균과 분위로 credible interval $[\,\mathrm{lo}=q_{0.05},\ \mathrm{hi}=q_{0.95}\,]$ 을 낸다.
(선택) 물리 가중 불확실성은 로그정규 전역인자 $w^{(r)}\!\sim\!\mathrm{LogN}$(평균 1, 변동계수
$cv$)로 합성한다.

\subsection{BYM2 공간 CAR (INLA식 Laplace 근사)}
\src{posterior.bym2\_spatial\_posterior}

켤레 사후는 이웃에서 정보를 빌리지 않는다. 이웃 격자 강도를 빌리도록(borrow strength)
공간 CAR 로 확장한다. 격자 $i$:
\begin{equation}
y_i\sim\mathrm{Poisson}\!\big(E_i\,e^{x_i}\big),\quad E_i=n_i\cdot(\text{전역율}),
\end{equation}
잠재 로그상대위험 $x$ 의 BYM2 prior 정밀도는 구조(ICAR)와 비구조(iid)의 보간:
\begin{equation}
Q_x(\tau,\phi)=\tau\big[(1-\phi)\,I+\phi\,Q_{\text{ICAR}}\big],\quad
Q_{\text{ICAR}}=D-W+\text{jitter}\cdot I.
\end{equation}
각 하이퍼 $(\tau,\phi)$ 에서 음의 로그사후 $f(x)=\sum_i(E_i e^{x_i}-y_i x_i)+\tfrac12 x^\top Q_x x$ 를
sparse Newton 으로 최소화해 모드 $x^\star$ 와 Hessian $H=\mathrm{diag}(E e^{x})+Q_x$ 를 얻고,
Laplace 주변우도
\begin{equation}
\log p(y\,|\,\tau,\phi)\approx \underbrace{\textstyle\sum_i(y_i x_i^\star-E_i e^{x_i^\star})}_{\text{Poisson loglik}}
+\tfrac12\log|Q_x|-\tfrac12 x^{\star\top}Q_x x^\star-\tfrac12\log|H|
\end{equation}
가 최대인 점을 채택한다. 격자 사후는 가우시안 $x_i\sim\mathcal N(x_i^\star,(H^{-1})_{ii})$,
상대배율은 로그정규 $e^{x_i}$.

% =====================================================================
\section{임계값·보정·평가}
% =====================================================================
\src{pfire/calibrate.py}

presence-only 라 모델은 \emph{순위}만 내고 $0/1$ 컷은 별도로 정한다.

\textbf{recall@top-$k$.} 위험 상위 $k\%$ 가 발화점-전주를 포함하는 비율(분리력 지표).

\textbf{단조 보정(PAVA isotonic).} 발화점=1, 표본 음성=0 으로 두고 위험점수$\to$확률을
단조(비감소) 회귀로 보정한다(절대확률 아닌 단조 보정값).

\textbf{F1 민감도.} 정답 양성비율 $\rho$ 미지 $\Rightarrow$ 가정별 곡선을 보고한다. 상위
$\rho$ 분위를 양성으로 두고
\begin{equation}
\text{prec}_{\text{proxy}}=\frac{\rho\cdot\text{recall}\cdot N}{n_{\text{pred}}},\quad
F_1=\frac{2\,\text{prec}_{\text{proxy}}\,\text{recall}}{\text{prec}_{\text{proxy}}+\text{recall}},
\quad \rho\in\{0.5,1,2,3,5,10\}\%.
\end{equation}

\textbf{per-regime conformal.} 교환가능 가정만으로 유한표본 커버리지를 보장한다. 체제 $r$
보정집합(=train fold 발화점에 매핑된 전주)의 위험점수를 오름차순 정렬해
\begin{equation}
\tau_r=\mathrm{score}_{(\lfloor\alpha(n_r+1)\rfloor)},\qquad
\Pr\!\big(\text{새 발화점 위험}\ge\tau_r\big)\gtrsim 1-\alpha,
\end{equation}
$\{p:\ \mathrm{risk}(p)\ge\tau_r\}$ 를 위험집합으로 둔다($\alpha=0.10\Rightarrow$ 명목 90\%
커버리지). 실측: 전체 0.902, 영동 0.907, 영서 0.902, 산간 0.897.

% =====================================================================
\section{상수 요약}
% =====================================================================
\src{pfire/config.py}

\begin{tabular}{lll}
\toprule
파라미터 & 값 & 의미 \\
\midrule
거리 감쇠 scale & 100/300/1500/3000\,m & 산림/도로/송전선/변전소 \\
게이트 온도 $\tau$ & 0.6 & 작을수록 하드 게이트 \\
$W$ 블렌드 $w_s$ & 0.82 & 시즌극값 : 일별동역학 \\
ISI/FFMC/양간 성분 & 0.05 / 0.75 / 0.20 & 일별 $W$ 결합 \\
LWR 상한 & 8.0 & Anderson 길이/폭비 cap \\
$W_\perp,\ L_{\text{back}}$ & 0.35 / 0.35\,km & 횡풍/풍상 스케일 \\
국지정규화 셀 & 4.0\,km & 영동 배경 제거 \\
노출 결합 $w$ & 0.25 & $R'=1-(1-R)(1-w\,e)$ \\
확산 컷오프 & 5.0\,km & 2019 고성 스케일 \\
conformal $\alpha$ & 0.10 & 명목 90\% 커버리지 \\
BYM2 $(\tau,\phi)$ 격자 & $\{1{\dots}300\}\times\{0.25{\dots}0.9\}$ & 정밀도·구조비 탐색 \\
SEED & 20260618 & 전 과정 결정성 \\
\bottomrule
\end{tabular}

\vspace{8pt}\hrule\vspace{4pt}
{\footnotesize 모든 식은 \texttt{pfire/} 모듈 구현과 일치하며, 발화점은 어디서도
per-pole 학습 라벨이 아니라 \emph{지역 집계 앵커·검증·임계값}으로만 쓰인다(누수 안전).}

\end{document}
